\chapter{拓扑空间与描述集合论初步}

\label{chap:01}

第 \ref{chap:00} 章已经把问题摆在了桌面上：一旦变量是函数、测度或轨道，而不再只是有限维向量，Boyd 母语里那些默认成立的事实就开始松动。最先失效的通常不是凸性本身，而是极限语言。我们在有限维里习惯于同时依赖序列、闭球、连续映射这些欧氏空间直觉；但在弱拓扑、弱$^\ast$拓扑和路径空间中，序列不再足够，这些话往往都需要重写。

因此，本章不是抽象的背景铺垫，而是在回答一个非常具体的问题：当优化变量已经离开 $\mathbb R^n$ 以后，我们究竟该用什么语言来谈闭包、连续性、紧性、可测性和“典型参数”？第一个答案是网与滤子，它们替代了“只谈序列”的习惯；第二个答案是 Baire 范畴，它解释为什么完备性会决定存在性和稳定性论证；第三个答案是 Polish、Lusin 与 Suslin 空间，它们给后续测度、随机过程和控制问题提供了真正稳定的状态空间舞台。

本章的推进仍然分成四步，但读的时候最好始终记着后文的用途。第 \ref{sec:1-1} 节解决“为什么弱拓扑里不能只谈序列”；第 \ref{sec:1-2} 节解决“为什么存在性论证会反复需要完备性和 Baire 方法”；第 \ref{sec:1-3} 节解释“为什么优化与随机过程的默认舞台往往是 Polish 空间”；第 \ref{sec:1-4} 节再说明“为什么还需要比 Polish 更宽的可测空间语言”。后文关于极小值存在、对偶构造、随机控制和图模型的很多接口，都会回到这一章来寻找最低层的支点。

\section{拓扑基础与网的收敛}

\label{sec:1-1}

第 \ref{chap:00} 章提出的第一个失效点，就是“弱拓扑里不能只谈序列”。这是无穷维分析里最容易被低估的一步：在欧氏空间中，闭包、连续性、紧性和聚点几乎都能靠序列表达，于是我们很自然地把“极限”理解成 $n\to\infty$。但只要进入一般拓扑空间，特别是后文会频繁出现的弱拓扑与弱$^\ast$拓扑，这种理解就不再够用。

因此，本节的核心不是再发明一种比序列更复杂的语言，而是把极限过程恢复到真正与拓扑相容的层面。网和滤子之所以重要，不是因为它们更抽象，而是因为它们能准确回答：当变量已经是函数、测度或轨道时，闭包究竟该怎样刻画，连续性究竟该怎样检查，紧性又究竟意味着什么。

\subsection{拓扑语言与定向集}

设 $X$ 为拓扑空间，$x\in X$。记 $\mathcal N(x)$ 为 $x$ 的邻域系，$\overline{A}$ 为子集 $A\subseteq X$ 的闭包。若对任意邻域 $U\in\mathcal N(x)$ 都有 $U\cap A\neq\varnothing$，则 $x\in \overline{A}$。这一定义不依赖于度量，因此比序列极限更基础。

\begin{definition}[定向集]\label{def:directed-set}
设 $D$ 为非空集合，$\preceq$ 为其上的预序关系，即满足自反性和传递性。若对任意 $\alpha,\beta\in D$ 都存在 $\gamma\in D$ 使得 $\alpha\preceq\gamma$ 且 $\beta\preceq\gamma$，则称 $(D,\preceq)$ 为一个\emph{定向集}。
\end{definition}

定向集的角色，是提供“最终足够大”这一概念，而不预设索引必须线性排列。自然数集 $(\mathbb N,\le)$ 是最基本的例子；更能反映一般拓扑结构的例子是邻域系 $\mathcal N(x)$ 配以逆包含顺序 $U\preceq V \iff V\subseteq U$。

\begin{definition}[网与网收敛]\label{def:net-convergence}
设 $(D,\preceq)$ 为定向集，$X$ 为拓扑空间。映射
\[
x_\bullet\colon D\to X,\qquad \alpha\mapsto x_\alpha
\]
称为 $X$ 中的一个\emph{网}。若对每个邻域 $U\in\mathcal N(x)$，都存在 $\alpha_0\in D$，使得当 $\alpha\succeq\alpha_0$ 时恒有 $x_\alpha\in U$，则称网 $(x_\alpha)_{\alpha\in D}$ \emph{收敛}到 $x$，记作 $x_\alpha\to x$。
\end{definition}

\begin{definition}[子网与滤子]\label{def:subnet-filter}
设 $(x_\alpha)_{\alpha\in D}$ 为网。若 $(E,\le)$ 为定向集，$\varphi\colon E\to D$ 满足
\begin{enumerate}
  \item $\varphi$ 保序，即 $\beta_1\le \beta_2$ 蕴含 $\varphi(\beta_1)\preceq \varphi(\beta_2)$；
  \item $\varphi(E)$ 在 $D$ 中共终，即对每个 $\alpha_0\in D$，都存在 $\beta_0\in E$ 使 $\varphi(\beta)\succeq \alpha_0$ 当 $\beta\succeq\beta_0$ 时成立，
\end{enumerate}
则称 $(x_{\varphi(\beta)})_{\beta\in E}$ 为原网的一个\emph{子网}。

另一方面，设 $X$ 为非空集合。若非空集合族 $\mathcal B\subseteq \mathcal P(X)$ 满足
\begin{enumerate}
  \item $\varnothing\notin \mathcal B$；
  \item 对任意 $B_1,B_2\in\mathcal B$，都存在 $B_3\in\mathcal B$ 使 $B_3\subseteq B_1\cap B_2$，
\end{enumerate}
则称 $\mathcal B$ 为 $X$ 上的一个\emph{滤子基}。滤子基的作用，是只记录“足够大的集合应该包含什么”，而不要求对上超集封闭。

若集合族 $\mathcal F\subseteq \mathcal P(X)$ 进一步满足
\begin{enumerate}
  \item $\varnothing\notin \mathcal F$；
  \item 若 $A,B\in\mathcal F$，则 $A\cap B\in\mathcal F$；
  \item 若 $A\in\mathcal F$ 且 $A\subseteq C\subseteq X$，则 $C\in\mathcal F$，
\end{enumerate}
则称 $\mathcal F$ 为 $X$ 上的一个\emph{滤子}。由任意滤子基 $\mathcal B$ 都可生成一个最小滤子
\[
\langle \mathcal B\rangle
:=
\{A\subseteq X:\exists B\in\mathcal B\text{ 使 }B\subseteq A\},
\]
称为由 $\mathcal B$ \emph{生成的滤子}。

在拓扑空间中，点 $x$ 的邻域系 $\mathcal N(x)$ 本身就是一个滤子；如果只取某个邻域基，则通常先得到的是滤子基。类似地，网的尾集族
\[
E_\alpha:=\{x_\beta:\beta\succeq\alpha\},\qquad \alpha\in D,
\]
正是一个典型的滤子基，因为定向性保证任意两个尾集都含有更小的公共尾集。由这些尾集生成的滤子，称为该网的\emph{尾滤子}。网与滤子在收敛论中携带的是同一类信息：一个网收敛到 $x$，当且仅当它的尾滤子收敛到邻域滤子 $\mathcal N(x)$；一个点是网的聚点，当且仅当某个比尾滤子更细的滤子收敛到该点。

本书在后续叙述中优先使用网，而不是直接使用滤子，主要有两个原因。第一，子网比“细化滤子”更适合写紧性和提取极限点的过程；第二，优化中的算法迭代、逼近族和约束逼近通常天然以索引族出现，用网语言书写更直接。不过一旦进入可测结构、完备化或抽象收敛论，滤子视角会帮助我们更清楚地识别“最终行为”究竟依赖于哪些集合论信息。
\end{definition}

可以把滤子理解为“忽略有限初段后剩下的全部信息”。对序列而言，尾滤子正是由集合
\[
\{x_n:n\ge N\},\qquad N\in\mathbb N,
\]
生成的；对邻域收敛而言，邻域滤子则记录了“一个点附近所有足够小的局部区域”。这两类滤子在形式上十分相似，这也是为什么网收敛与邻域结构能够自然拼接在一起。

\subsection{闭包与连续性的网刻画}

\begin{proposition}[闭包的网刻画]\label{prop:closure-net}
设 $A\subseteq X$，$x\in X$。则
\[
x\in\overline{A}
\quad\Longleftrightarrow\quad
\text{存在网 }(x_\alpha)_{\alpha\in D}\subseteq A\text{ 使 }x_\alpha\to x.
\]
\end{proposition}

\begin{proof}
若存在这样的网，则对任意邻域 $U\in\mathcal N(x)$，收敛性保证存在 $\alpha_0$ 使得 $\alpha\succeq\alpha_0$ 时 $x_\alpha\in U\cap A$，从而 $U\cap A\neq\varnothing$，故 $x\in\overline{A}$。

反之，若 $x\in\overline{A}$，则对每个邻域 $U\in\mathcal N(x)$ 都可选取一点 $x_U\in U\cap A$。将 $\mathcal N(x)$ 赋予逆包含顺序：$U\preceq V$ 当且仅当 $V\subseteq U$。这使 $\mathcal N(x)$ 成为定向集，而 $(x_U)_{U\in\mathcal N(x)}$ 是 $A$ 中的网。对任意邻域 $W\in\mathcal N(x)$，若 $U\preceq W$，即 $U\subseteq W$，则 $x_U\in U\subseteq W$。故该网收敛到 $x$。
\end{proof}

\begin{proposition}[连续性的网刻画]\label{prop:continuity-net}
设 $f\colon X\to Y$ 为映射，$x\in X$。则 $f$ 在 $x$ 处连续，当且仅当对任意满足 $x_\alpha\to x$ 的网 $(x_\alpha)$，都有 $f(x_\alpha)\to f(x)$。
\end{proposition}

\begin{proof}
先设 $f$ 在 $x$ 处连续。任取 $f(x)$ 的邻域 $V$，由连续性可取 $x$ 的邻域 $U$ 使 $f(U)\subseteq V$。若 $x_\alpha\to x$，则存在 $\alpha_0$ 使得 $\alpha\succeq\alpha_0$ 时 $x_\alpha\in U$，于是 $f(x_\alpha)\in V$。故 $f(x_\alpha)\to f(x)$。

反之，假设像网保持收敛。若 $f$ 在 $x$ 不连续，则存在 $f(x)$ 的邻域 $V$ 使得对任意 $x$ 的邻域 $U$ 都有 $U\not\subseteq f^{-1}(V)$。于是每个 $U\in\mathcal N(x)$ 中都可选取 $x_U\in U\setminus f^{-1}(V)$。与命题~\ref{prop:closure-net} 证明中相同，$(x_U)_{U\in\mathcal N(x)}$ 收敛到 $x$，但其像网始终落在 $Y\setminus V$ 中，不可能收敛到 $f(x)$。矛盾。
\end{proof}

\begin{proposition}[聚点与子网]\label{prop:cluster-subnet}
设 $(x_\alpha)_{\alpha\in D}$ 为 $X$ 中的网，$x\in X$。则下列命题等价：
\begin{enumerate}
  \item 对任意邻域 $U\in\mathcal N(x)$ 与任意 $\alpha_0\in D$，都存在 $\alpha\succeq\alpha_0$ 使 $x_\alpha\in U$；
  \item 原网存在一个子网收敛到 $x$。
\end{enumerate}
\end{proposition}

\begin{proof}
$(2)\Rightarrow(1)$ 显然，因为收敛子网的尾部落在每个邻域中，而子网索引映射共终。

下证 $(1)\Rightarrow(2)$。令
\[
E=\{(\alpha,U):\alpha\in D,\ U\in\mathcal N(x),\ x_\alpha\in U\},
\]
并定义
\[
(\alpha_1,U_1)\preceq(\alpha_2,U_2)
\quad\Longleftrightarrow\quad
\alpha_1\preceq \alpha_2\ \text{且}\ U_2\subseteq U_1.
\]
由条件 (1) 可知 $E$ 非空，且对任意两个元素都可找到共同上界，所以 $E$ 是定向集。定义 $\varphi\colon E\to D$ 为 $\varphi(\alpha,U)=\alpha$，则 $\varphi$ 保序且共终，因此 $(x_{\varphi(\beta)})_{\beta\in E}$ 是原网的子网。对任意邻域 $W\in\mathcal N(x)$，若 $\beta=(\alpha,U)\succeq(\alpha_0,W)$，则必有 $U\subseteq W$，因而 $x_{\varphi(\beta)}=x_\alpha\in U\subseteq W$。故此子网收敛到 $x$。
\end{proof}

\begin{proposition}[第一可数空间中序列已足够]\label{prop:first-countable-sequence}
若 $X$ 是第一可数空间，$A\subseteq X$，则对任意 $x\in X$，
\[
x\in \overline{A}
\quad\Longleftrightarrow\quad
\text{存在序列 }(a_n)_{n\ge1}\subseteq A\text{ 使 }a_n\to x.
\]
\end{proposition}

\begin{proof}
由命题~\ref{prop:closure-net} 只需证明右向左的逆向构造。若 $x\in\overline{A}$，取 $x$ 的可数邻域基 $(U_n)_{n\ge1}$，并令
\[
V_n:=U_1\cap\cdots\cap U_n.
\]
则每个 $V_n$ 仍是 $x$ 的邻域，且 $V_{n+1}\subseteq V_n$。由 $x\in\overline{A}$ 可取 $a_n\in A\cap V_n$。于是对任意邻域 $U$，存在某个 $N$ 使 $U_N\subseteq U$，从而当 $n\ge N$ 时 $a_n\in V_n\subseteq U_N\subseteq U$。故 $a_n\to x$。
\end{proof}

\paragraph{为何弱拓扑中序列语言不够}

\begin{remark}
命题~\ref{prop:first-countable-sequence} 表明，序列之所以在欧氏空间中万能，并不是因为“极限天然应该由序列描述”，而是因为这些空间具备第一可数性。后文讨论弱拓扑与弱$^\ast$拓扑时，这一性质往往失效。例如第 \ref{sec:3-2} 节定理~\ref{thm:banach-alaoglu} 给出的弱$^\ast$紧性，在非可分情形通常只能保证网有收敛子网，而不能保证任意序列都有收敛子列。因此本书后文凡涉及弱收敛、弱下半连续性与对偶紧性，默认都应使用网语言，而不是未经说明地退回序列语言。
\end{remark}

\begin{example}[最小抽象例子：弱$^\ast$极限网]\label{ex:weak-star-net}
设 $X$ 为赋范空间，$B_{X^\ast}$ 为对偶空间单位球。第 \ref{sec:3-2} 节定理~\ref{thm:banach-alaoglu} 说明 $B_{X^\ast}$ 在弱$^\ast$拓扑下紧。这里“紧”的自然含义是：任意网都存在收敛子网。把这个例子放在本节，是为了提醒读者后文出现的极小化列、乘子列或近似解族，若只在弱$^\ast$拓扑下有紧性，就不能只问“是否存在收敛子列”，而必须问“是否存在收敛子网”。
\end{example}

\begin{example}[有限维坍缩：欧氏空间中的网与序列]\label{ex:euclidean-net}
在 $\mathbb R^n$ 的通常拓扑中，第一可数性成立，因此命题~\ref{prop:first-countable-sequence} 告诉我们：闭包、连续性与聚点完全可以由序列刻画。把这个例子放在这里，是为了明确后文与 Boyd 书相接时，为什么很多有限维论证看上去只需处理序列；那并非一般原理，而是欧氏空间的特殊福利。
\end{example}

\begin{example}[算法触点：函数空间中的迭代极限]\label{ex:function-space-iteration}
在 Hilbert 空间上研究投影算法或近端算法时，迭代点常常只能先证明有界，再借助弱拓扑获得极限点。若空间或拓扑选择得更一般，例如考虑测度空间上的概率流或策略空间上的松弛解，极限对象很可能首先通过网出现。把这个例子放在这里，是为了给第 9 章与第 10 章的算子分裂算法预留正确的收敛语言接口。
\end{example}

到这里，本节其实已经回答了第 \ref{chap:00} 章中的第一个问题：为什么一旦离开有限维，连“收敛”都必须重写。下一节将继续回答第二个问题：即便有了正确的极限语言，为什么存在性和稳定性论证仍然离不开完备性与 Baire 方法。

\subsection*{有限维对照}

Boyd 书中关于闭集、极限点和连续映射的讨论，默认都在 $\mathbb R^n$ 这一第一可数环境中完成。因此其结论可以直接用序列表述而不损失信息。本节的作用不是推翻这些有限维叙述，而是指出它们在无穷维中的真实上位版本：定义~\ref{def:net-convergence}、命题~\ref{prop:closure-net} 和命题~\ref{prop:continuity-net}。

\subsection*{习题（待补）}

\begin{exercise}[闭包与子网的练习占位]\label{exr:closure-subnet-placeholder}
补一题：要求证明命题~\ref{prop:cluster-subnet} 在序列情形退化为“聚点当且仅当存在收敛子列”，并指出证明中哪些地方依赖第一可数性。
\end{exercise}

\begin{exercise}[弱拓扑桥接题占位]\label{exr:weak-topology-placeholder}
补一题：把命题~\ref{prop:continuity-net} 落回欧氏空间，比较“网连续”“序列连续”和通常 $\varepsilon$-$\delta$ 连续性的等价性。
\end{exercise}


\section{Baire 纲定理及其推论}

\label{sec:1-2}

上一节解决了“极限该怎么说”，但这还不够。第 \ref{chap:00} 章里的第二个失效点是存在性：即便我们已经知道弱极限、子网和闭包该怎样表述，仍然需要某种结构来保证“坏集合不能处处出现”。这正是完备性和 Baire 范畴方法进入优化的地方。

Baire 方法之所以在泛函分析中占据中心地位，并不是因为它提供了某种旁门技巧，而是因为它把完备性转化成了一条结构结论：可数个过于稀薄的集合，不可能填满一个真正良好的空间。后文许多看似与“范畴”无关的结论，例如一致有界、开映射、闭图像以及典型参数上的稳定性，底层都在使用这一点。

\subsection{稀疏集、第一纲集与 Baire 空间}

\begin{definition}[Baire 范畴的基本概念]\label{def:baire-space}
设 $X$ 为拓扑空间，$A\subseteq X$。
\begin{enumerate}
  \item 若 $\operatorname{int}\overline{A}=\varnothing$，则称 $A$ 为\emph{无处稠密集}。
  \item 若 $A$ 可以写成可数个无处稠密集的并，则称 $A$ 为\emph{第一纲集}；否则称其为\emph{第二纲集}。
  \item 若 $R\subseteq X$ 的补集是第一纲集，则称 $R$ 为\emph{剩余集}。
  \item 若任意可数个稠密开集的交仍稠密，则称 $X$ 为\emph{Baire 空间}。
\end{enumerate}
\end{definition}

在 Baire 空间中，剩余集总是稠密的。因此“典型性质”常常被理解为“在一个剩余集上成立”。后文处理连续线性算子族、值函数正则性和参数稳定性时，这种语言会反复出现。

\subsection{完备度量空间中的 Baire 纲定理}

\begin{theorem}[完备度量空间上的 Baire 纲定理]\label{thm:baire-complete}
设 $(X,d)$ 为非空完备度量空间，$\{U_n\}_{n\ge 1}$ 为 $X$ 中一列稠密开集。则
\[
\bigcap_{n=1}^{\infty} U_n
\]
在 $X$ 中稠密。
\end{theorem}

\begin{proof}
任取非空开集 $V\subseteq X$。只需证明
\[
V\cap \bigcap_{n=1}^{\infty} U_n\neq \varnothing.
\]
因为 $U_1$ 稠密且开，故可在 $V\cap U_1$ 中选取闭球 $\overline{B}(x_1,r_1)$，满足
\[
\overline{B}(x_1,r_1)\subseteq V\cap U_1,\qquad r_1\le 2^{-1}.
\]
递归地，若已选得 $\overline{B}(x_n,r_n)\subseteq U_n\cap B(x_{n-1},r_{n-1})$ 且 $r_n\le 2^{-n}$，则由 $U_{n+1}$ 稠密开，可在 $U_{n+1}\cap B(x_n,r_n)$ 中选取闭球
\[
\overline{B}(x_{n+1},r_{n+1})\subseteq U_{n+1}\cap B(x_n,r_n),\qquad r_{n+1}\le 2^{-(n+1)}.
\]

于是得到一列递减的非空闭集
\[
F_n:=\overline{B}(x_n,r_n),\qquad F_{n+1}\subseteq F_n,\qquad \operatorname{diam}(F_n)\le 2r_n\to 0.
\]
任取 $y_n\in F_n$。对 $m>n$，有 $y_m,y_n\in F_n$，故
\[
d(y_m,y_n)\le \operatorname{diam}(F_n)\to 0.
\]
因此 $(y_n)$ 为 Cauchy 列，由完备性收敛到某个 $x\in X$。又因每个 $F_n$ 闭且对所有 $m\ge n$ 有 $y_m\in F_n$，故 $x\in F_n$ 对一切 $n$ 都成立。于是
\[
x\in F_1\subseteq V,\qquad x\in F_n\subseteq U_n\quad (n\ge 1),
\]
故 $x\in V\cap\bigcap_{n\ge 1}U_n$。
\end{proof}

\begin{example}[为什么完备性不能省略]\label{ex:baire-rationals}
在有理数空间 $(\mathbb Q,|\cdot|)$ 中，枚举 $\mathbb Q=\{q_n\}_{n\ge1}$ 并令 $U_n=\mathbb Q\setminus\{q_n\}$。每个 $U_n$ 都是稠密开集，但
\[
\bigcap_{n=1}^{\infty}U_n=\varnothing.
\]
把这个例子放在这里，是为了强调定理~\ref{thm:baire-complete} 的核心前提不是“可数交”，而是完备性本身。
\end{example}

\subsection{局部紧 Hausdorff 空间的版本}

\begin{lemma}[局部紧 Hausdorff 空间中的紧闭包开集]\label{lem:locally-compact-open}
设 $X$ 为局部紧 Hausdorff 空间，$O\subseteq X$ 为非空开集。则存在非空开集 $W$ 使得
\[
\overline{W}\subseteq O,\qquad \overline{W}\ \text{紧}.
\]
\end{lemma}

\begin{proof}
任取 $x\in O$。由局部紧性可取 $x$ 的开邻域 $V$ 使得 $\overline{V}$ 紧；由 Hausdorff 性，紧集是闭集。再以正则性选取开集 $W$ 满足
\[
x\in W\subseteq \overline{W}\subseteq V\cap O.
\]
于是 $\overline{W}$ 为紧集且包含于 $O$。
\end{proof}

\begin{theorem}[局部紧 Hausdorff 空间上的 Baire 纲定理]\label{thm:baire-locally-compact}
设 $X$ 为非空局部紧 Hausdorff 空间，$\{U_n\}_{n\ge1}$ 为 $X$ 中一列稠密开集。则 $\bigcap_{n\ge1}U_n$ 在 $X$ 中稠密。
\end{theorem}

\begin{proof}
任取非空开集 $V\subseteq X$。由引理~\ref{lem:locally-compact-open}，存在非空开集 $W_1$ 使
\[
\overline{W_1}\subseteq V\cap U_1,\qquad \overline{W_1}\text{ 紧}.
\]
递归地，若已构造出非空开集 $W_n$ 使 $\overline{W_n}$ 紧且 $\overline{W_n}\subseteq U_n\cap W_{n-1}$，则对非空开集 $W_n\cap U_{n+1}$ 再次应用引理~\ref{lem:locally-compact-open}，可得非空开集 $W_{n+1}$ 满足
\[
\overline{W_{n+1}}\subseteq W_n\cap U_{n+1},\qquad \overline{W_{n+1}}\text{ 紧}.
\]
于是 $\{\overline{W_n}\}_{n\ge1}$ 构成递减的非空紧集列。紧性保证有限交性质成立，因此
\[
\bigcap_{n=1}^{\infty}\overline{W_n}\neq\varnothing.
\]
取 $x$ 属于该交，则 $x\in \overline{W_1}\subseteq V$ 且对每个 $n$ 有 $x\in \overline{W_n}\subseteq U_n$。故
\[
x\in V\cap\bigcap_{n=1}^{\infty}U_n,
\]
从而交集稠密。
\end{proof}

\subsection{等价刻画与前瞻性推论}

\begin{proposition}[Baire 空间的等价刻画]\label{prop:baire-equivalent}
设 $X$ 为拓扑空间。下列命题等价：
\begin{enumerate}
  \item $X$ 是 Baire 空间；
  \item 任意第一纲集都没有内点；
  \item 任意剩余集都是稠密集；
  \item 任意非空开子集都是第二纲集。
\end{enumerate}
\end{proposition}

\begin{proof}
$(1)\Rightarrow(2)$：若 $M=\bigcup_{n\ge1}A_n$ 为第一纲集，其中每个 $A_n$ 无处稠密，则 $X\setminus \overline{A_n}$ 均为稠密开集。由 (1) 可知
\[
\bigcap_{n\ge1}(X\setminus\overline{A_n})
\]
稠密，于是其补集
\[
\bigcup_{n\ge1}\overline{A_n}
\]
没有内点。因为 $M\subseteq \bigcup_{n\ge1}\overline{A_n}$，故 $M$ 也没有内点。

$(2)\Rightarrow(3)$：若 $R$ 为剩余集，则 $X\setminus R$ 为第一纲集，由 (2) 知其内点为空，等价于 $R$ 稠密。

$(3)\Rightarrow(1)$：任意可数个稠密开集的交，其补为可数个闭无内点集之并，故是第一纲集。于是该交为剩余集，由 (3) 稠密。

$(2)\Rightarrow(4)$：若非空开集 $O$ 是第一纲集，则它作为自身的子集有非空内点，与 (2) 矛盾。

$(4)\Rightarrow(2)$：若第一纲集 $M$ 有非空内点 $O$，则 $O\subseteq M$，从而 $O$ 也是第一纲集，与 (4) 矛盾。
\end{proof}

\begin{corollary}[稠密 $G_\delta$ 集的典型性]\label{cor:dense-gdelta}
在 Baire 空间中，可数个稠密开集的交是稠密的 $G_\delta$ 集；因而任意剩余集都包含一个稠密 $G_\delta$ 子集。
\end{corollary}

\begin{proof}
前一结论来自定义；后一结论只需把剩余集写成第一纲集的补，再将对应的闭包补取交即可。
\end{proof}

\paragraph{对后文的前瞻}

\begin{remark}
开映射定理、闭图像定理与一致有界原理的共同底层机制，正是定理~\ref{thm:baire-complete} 与命题~\ref{prop:baire-equivalent}。本书在第 3 章讨论 Banach 空间工具箱时，将重新回到本节，不在这里提前展开线性空间版本的全部证明。
\end{remark}

\begin{example}[最小抽象例子：一致有界原理的拓扑影子]\label{ex:uniform-boundedness-shadow}
设 $\{T_i\}_{i\in I}$ 为 Banach 空间上的连续线性算子族。若对每个固定向量 $x$，集合 $\{\|T_i x\|:i\in I\}$ 有界，则“范数爆炸点集”可以写成可数个闭集的并。Baire 方法排除了这些闭集在整个空间中都太薄却又处处出现的可能，因此导出算子范数的一致有界性。把这个例子放在这里，是为了说明范畴方法不是抽象装饰，而是泛函分析主定理的发动机。
\end{example}

\begin{example}[有限维坍缩：欧氏空间中一切自动成为 Baire]\label{ex:finite-dimensional-baire}
任意欧氏空间都是完备度量空间，因此由定理~\ref{thm:baire-complete} 自动是 Baire 空间。把这个例子放在这里，是为了说明有限维线性代数中“满射线性映射自动开”“图像闭合行为良好”等现象，其实早已暗含在本节的完备性框架里，只是欧氏空间把这些前提隐藏得过于自然。
\end{example}

\begin{example}[应用触点：坏参数集很薄]\label{ex:thin-bad-parameters}
在优化模型的参数分析里，常见结论是“对典型参数，解映射满足某种正则性；例外参数至多构成第一纲集”。把这个例子放在这里，是为了给后文的对偶稳定性与值函数敏感性分析提供术语准备：所谓“典型”，往往并非指测度意义上的几乎处处，而是 Baire 范畴意义上的剩余集。
\end{example}

因此，本节回答的是第 \ref{chap:00} 章中的第二个问题：为什么存在性与稳定性论证需要某种“坏集合不能太大”的原理。再往前一步，优化和随机过程还要求空间既能承受完备性，又有足够好的可分性和 Borel 结构；这就把我们带到下一节的 Polish 空间。

\subsection*{有限维对照}

Boyd 书几乎不需要显式提到 Baire 范畴，因为有限维范数空间自动完备，且很多线性映射性质可由矩阵论直接验证。本节的价值在于指出：这些熟悉结论在无穷维里不再是线性代数常识，而是定理~\ref{thm:baire-complete} 与定理~\ref{thm:baire-locally-compact} 的推论。

\subsection*{习题（待补）}

\begin{exercise}[Baire 等价刻画占位]\label{exr:baire-equivalence-placeholder}
补一题：要求把命题~\ref{prop:baire-equivalent} 的四个条件补成完整循环证明，并解释“剩余”与“第一纲”之间的对偶关系。
\end{exercise}

\begin{exercise}[完备性桥接题占位]\label{exr:baire-complete-placeholder}
补一题：比较 $\mathbb R$ 与 $\mathbb Q$ 中稠密开集可数交的行为，说明为什么例~\ref{ex:baire-rationals} 恰好刻画了完备性的缺失。
\end{exercise}


\section{完备性与 Polish 空间}

\label{sec:1-3}

到这里，读者已经看到：仅有“正确的极限语言”和“完备性带来的 Baire 性质”还不足以支撑后文全部内容。第 \ref{chap:00} 章真正要追问的是：为什么优化与随机过程的默认舞台常常不是任意 Banach 空间，而是 Polish 空间？原因很直接。后文既要讨论极小值存在，又要讨论 Borel 结构、概率测度和可测选择；这就要求状态空间同时具备完备性、可分性和稳定的可测结构。

Polish 空间正是这些要求的交汇点。它既保留了上一节的完备性力量，又给出了随机过程、路径空间和测度论真正可操作的状态空间类别。对本书来说，Polish 空间不是描述集合论中的漂亮名词，而是“把优化对象、概率对象和拓扑对象放在同一个舞台上”的最常用答案。

\subsection{可完备度量化与 Polish 空间}

\begin{definition}[可完备度量化]\label{def:completely-metrizable}
拓扑空间 $X$ 称为\emph{可完备度量化}的，若存在一个与其拓扑相容的完备度量 $d$。换言之，$X$ 的拓扑可以由某个完备度量诱导出来。
\end{definition}

\begin{definition}[Polish 空间]\label{def:polish-space}
若拓扑空间 $X$ 既可分又可完备度量化，则称 $X$ 为\emph{Polish 空间}。
\end{definition}

由定理~\ref{thm:baire-complete} 可知，每个可完备度量化空间都是 Baire 空间；因此每个 Polish 空间自动满足命题~\ref{prop:baire-equivalent} 的全部等价条件。换言之，Polish 空间不仅“测度上好用”，也“范畴上不坏”。

\subsection{开子空间与 $G_\delta$ 子集}

\begin{theorem}[开子空间的可完备度量化]\label{thm:open-subspace-completely-metrizable}
设 $X$ 为可完备度量化空间，$U\subseteq X$ 为开集。则 $U$ 也是可完备度量化空间。
\end{theorem}

\begin{proof}
取与 $X$ 的拓扑相容的完备度量 $d$。若 $U=X$ 结论显然，下面设 $U\neq X$，并记 $F:=X\setminus U$。对 $x\in U$ 定义
\[
\Phi(x)=\left(x,\frac{1}{d(x,F)}\right)\in X\times \mathbb R.
\]
由于 $F$ 闭，函数 $x\mapsto d(x,F)$ 在 $U$ 上连续且严格为正，因此 $\Phi$ 连续。显然它是单射，而投影到第一坐标给出其像上的连续逆映射，所以 $\Phi$ 是 $U$ 到 $\Phi(U)$ 的同胚。

下证 $\Phi(U)$ 在 $X\times\mathbb R$ 中闭。设 $\Phi(x_n)\to (x,t)$。则 $x_n\to x$ 且 $1/d(x_n,F)\to t$。由于距离到闭集的函数连续，得到
\[
d(x_n,F)\to d(x,F).
\]
若 $x\in F$，则 $d(x,F)=0$，从而 $d(x_n,F)\to 0$，这将迫使 $1/d(x_n,F)\to +\infty$，与其在 $\mathbb R$ 中收敛到 $t$ 矛盾。因此 $d(x,F)>0$，即 $x\in U$。再由连续性可得
\[
t=\lim_{n\to\infty}\frac{1}{d(x_n,F)}=\frac{1}{d(x,F)}.
\]
所以 $(x,t)=\Phi(x)\in \Phi(U)$，闭性得证。

因为 $X\times\mathbb R$ 在乘积拓扑下仍可由完备度量诱导，而闭子空间保持完备性，故 $\Phi(U)$ 可完备度量化。再由 $U\cong \Phi(U)$，得到 $U$ 可完备度量化。
\end{proof}

\begin{corollary}[开子空间的 Polish 性]\label{cor:open-polish}
Polish 空间的任意开子空间仍是 Polish 空间。
\end{corollary}

\begin{proof}
开子空间继承可分性；再由定理~\ref{thm:open-subspace-completely-metrizable} 继承可完备度量化。
\end{proof}

\begin{proposition}[可数乘积保持 Polish 性]\label{prop:countable-product-polish}
设 $\{X_n\}_{n\ge1}$ 为一列 Polish 空间，则乘积空间 $\prod_{n\ge1}X_n$ 仍是 Polish 空间。
\end{proposition}

\begin{proof}
对每个 $n$ 取有界完备度量 $d_n\le 1$ 诱导 $X_n$ 的拓扑。定义
\[
D(x,y)=\sum_{n=1}^{\infty}2^{-n}d_n(x_n,y_n),\qquad x=(x_n),\ y=(y_n).
\]
则 $D$ 诱导乘积拓扑。若 $(x^{(m)})_{m\ge1}$ 是 $D$-Cauchy 列，则对每个固定 $n$，坐标列 $x_n^{(m)}$ 在 $(X_n,d_n)$ 中为 Cauchy，从而收敛到某个 $x_n\in X_n$。由逐坐标收敛和定义可知 $x^{(m)}\to (x_n)$，故 $D$ 完备。可分性来自每个 $X_n$ 的可数稠密子集取有限支近似的标准构造。
\end{proof}

\begin{theorem}[$G_\delta$ 子集的 Polish 性]\label{thm:gdelta-polish}
设 $X$ 为 Polish 空间，$G\subseteq X$ 为 $G_\delta$ 子集。则 $G$ 仍为 Polish 空间。
\end{theorem}

\begin{proof}
取开集列 $\{U_n\}_{n\ge1}$ 使
\[
G=\bigcap_{n=1}^{\infty}U_n.
\]
由推论~\ref{cor:open-polish}，每个 $U_n$ 都是 Polish 空间。考虑对角嵌入
\[
\Delta\colon G\to \prod_{n=1}^{\infty}U_n,\qquad \Delta(x)=(x,x,\dots).
\]
这是一个同胚到其像的映射。其像恰是满足各坐标相等的对角集；由于乘积空间是 Hausdorff 的，对角集为闭集。由命题~\ref{prop:countable-product-polish}，$\prod_n U_n$ 是 Polish 空间，闭子空间仍是 Polish 空间，因此 $G$ 是 Polish 空间。
\end{proof}

\paragraph{关于 $G_\delta$ 判别}

\begin{remark}
定理~\ref{thm:gdelta-polish} 给出了本书最常用的一个判别原则：当一个空间已经嵌在某个 Polish 空间里时，只要能证明它是一个 $G_\delta$ 子集，往往就已经足够恢复完整的 Polish 结构。后文讨论可行集、路径空间和策略空间时，这条原则会反复使用。
\end{remark}

\subsection{典型模型}

\begin{example}[欧氏空间与可分 Banach 空间]\label{ex:rn-separable-banach}
$\mathbb R^n$ 在通常欧氏度量下是完备且可分的，因此是 Polish 空间。更一般地，任何可分 Banach 空间在其范数度量下也是 Polish 空间。把这些例子放在这里，是为了说明 Boyd 书的整个有限维世界，以及后文大部分 Hilbert 空间算法模型，天然都生活在定义~\ref{def:polish-space} 的框架之中。
\end{example}

\begin{example}[标准 Baire 空间 $\mathbb N^{\mathbb N}$]\label{ex:baire-space}
将 $\mathbb N$ 赋予离散拓扑，并在乘积空间 $\mathbb N^{\mathbb N}$ 上考虑度量
\[
d(\alpha,\beta)=\sum_{n=1}^{\infty}2^{-n}\mathbf 1_{\{\alpha_n\neq \beta_n\}}.
\]
则这是一个完备可分度量空间，因此是 Polish 空间。把它放在这里，是为了给下一节做准备：标准 Baire 空间在描述集合论中扮演“可数编码空间”的角色，很多 Suslin 与 Lusin 结构都可以从它出发进行参数化。
\end{example}

\begin{example}[应用触点：路径空间与随机优化]\label{ex:path-space-polish}
设 $C([0,T])$ 配以一致范数，则它是可分 Banach 空间，故由例~\ref{ex:rn-separable-banach} 为 Polish 空间。把这个例子放在这里，是为了说明后文最优控制与随机过程的状态轨道空间，往往并不是抽象的“任意函数集”，而是一个具有稳定 Borel 结构的 Polish 空间。
\end{example}

本节因此回答了第 \ref{chap:00} 章中的第三个问题：为什么后文在谈随机过程、控制轨道和可测优化时，会反复把 Polish 空间当作默认舞台。下一节还要继续向外推广一步，因为实际建模里出现的状态空间，未必总是以 Polish 形式直接给出，它们常常更像某个 Polish 空间的连续像。

\subsection*{有限维对照}

在有限维中，“开集仍然好”“可数交仍然可控”这些性质几乎从不被单独强调，因为 $\mathbb R^n$ 的闭子集、开子集和很多自然出现的约束集都自动继承良好结构。本节把这种直觉抽象成三个真正可复用的结论：定理~\ref{thm:open-subspace-completely-metrizable}、命题~\ref{prop:countable-product-polish} 与定理~\ref{thm:gdelta-polish}。

\subsection*{习题（待补）}

\begin{exercise}[开子空间完备化占位]\label{exr:open-subspace-placeholder}
补一题：要求用定理~\ref{thm:open-subspace-completely-metrizable} 构造 $(0,1)$ 上一个与通常拓扑相容但完备的度量。
\end{exercise}

\begin{exercise}[$G_\delta$ 判别占位]\label{exr:gdelta-placeholder}
补一题：证明某个具体函数约束集是 Banach 空间中的 $G_\delta$ 子集，并据此应用定理~\ref{thm:gdelta-polish} 得到其 Polish 性。
\end{exercise}


\section{Baire 空间与 Lusin 空间}

\label{sec:1-4}

Polish 空间已经足以覆盖很多优化与概率问题，但第 \ref{chap:00} 章中的最后一个问题还没有完全回答：为什么后文有时还需要比 Polish 更宽的可测空间语言？原因在于，真实建模里出现的状态空间并不总是直接以“可分完备度量空间”出现。它们更常见的面貌，是某个良好空间经过连续映射、参数化或编码之后得到的像。

因此，本节并不想把描述集合论再铺成一整门课程，而是只建立一条对后文真正有用的层级：
\[
\text{Polish} \Longrightarrow \text{Lusin} \Longrightarrow \text{Suslin}.
\]
这条层级的作用，是在尽量保留 Polish 空间可测正则性的同时，允许状态空间经过更现实的连续像变形。后文讨论随机控制、可测选择和参数化策略空间时，真正需要的往往正是这种“比 Polish 更宽，但还没有宽到失控”的环境。

\subsection{标准 Baire 空间的角色}

\begin{definition}[标准 Baire 空间]\label{def:standard-baire-space}
记
\[
\mathcal N:=\mathbb N^{\mathbb N}
\]
并赋予其乘积拓扑，则称 $\mathcal N$ 为\emph{标准 Baire 空间}。由例~\ref{ex:baire-space}，它是一个 Polish 空间。
\end{definition}

标准 Baire 空间的重要性在于：它提供了一个统一的可数编码平台。后文若要把控制策略、可测选择或参数族写成“由一串整数决定的对象”，实质上就是把问题转移到 $\mathcal N$ 上。

\begin{proposition}[标准 Baire 空间的满射性]\label{prop:baire-space-surjection}
每个非空 Polish 空间都是标准 Baire 空间的连续像。
\end{proposition}

\begin{proof}
证明可由“逐层缩小的可数开覆盖”完成：先对给定 Polish 空间 $X$ 取一族直径趋于零的可数开覆盖，再沿有限自然数列递归地选取嵌套开集 $U_s$，使得
\[
U_s=\bigcup_{n\ge1}U_{s^\frown n},\qquad \operatorname{diam}(U_s)\to 0 \text{ 当 }|s|\to\infty.
\]
对每个 $\alpha=(\alpha_1,\alpha_2,\dots)\in\mathbb N^{\mathbb N}$，完备性保证
\[
\bigcap_{k\ge1}\overline{U_{\alpha|k}}
\]
恰含一点，据此定义映射 $\pi(\alpha)$。直径控制给出唯一性与连续性，覆盖性则保证满射。这里不展开全部技术细节，因为后文只需要这一结论作为参数化工具，而不需要其完整构造。
\end{proof}

\subsection{Suslin 空间作为过渡层}

\begin{definition}[Suslin 空间]\label{def:suslin-space}
拓扑空间 $X$ 称为\emph{Suslin 空间}，若存在 Polish 空间 $P$ 及连续满射 $f\colon P\to X$。
\end{definition}

Suslin 空间的关键词是“连续像”。这一定义非常宽松，却保留了大量对测度论有用的正则性；尤其是连续像这一操作，在建模时远比“是某个 Banach 空间的闭子集”更常见。

\begin{definition}[Lusin 空间]\label{def:lusin-space}
Hausdorff 空间 $X$ 称为\emph{Lusin 空间}，若存在 Polish 空间 $P$ 及连续双射 $f\colon P\to X$。
\end{definition}

Lusin 空间介于 Polish 与 Suslin 之间。与 Suslin 空间相比，它保留了“来自 Polish 空间的点级编码几乎没有丢失”的特征；与 Polish 空间相比，它允许拓扑结构经过较温和的变形。

\begin{proposition}[层级关系：Polish、Lusin 与 Suslin]\label{prop:polish-lusin-suslin}
任意 Polish 空间都是 Lusin 空间，任意 Lusin 空间都是 Suslin 空间。
\end{proposition}

\begin{proof}
若 $X$ 是 Polish 空间，则取恒等映射 $\mathrm{id}_X\colon X\to X$，可见 $X$ 满足定义~\ref{def:lusin-space}，故为 Lusin 空间。若 $X$ 是 Lusin 空间，则按定义存在 Polish 空间 $P$ 与连续双射 $f\colon P\to X$。双射当然是满射，因此 $X$ 也满足定义~\ref{def:suslin-space}，故为 Suslin 空间。
\end{proof}

\begin{proposition}[Suslin 与 Lusin 空间的基本保持性]\label{prop:suslin-lusin-stability}
\begin{enumerate}
  \item Suslin 空间的连续像仍为 Suslin 空间。
  \item 若 $X$ 是 Lusin 空间，则 $X$ 的任意开子空间、闭子空间及 $G_\delta$ 子空间仍是 Lusin 空间。
\end{enumerate}
\end{proposition}

\begin{proof}
对 (1)，设 $X$ 是 Suslin 空间，$f\colon P\to X$ 为 Polish 空间 $P$ 上的连续满射，且 $g\colon X\to Y$ 连续。则 $g\circ f\colon P\to Y$ 连续且满射到 $g(X)$，故 $g(X)$ 为 Suslin 空间。

对 (2)，设 $X$ 是 Lusin 空间，取 Polish 空间 $P$ 与连续双射 $f\colon P\to X$。若 $A\subseteq X$ 为开集，则 $f^{-1}(A)$ 在 $P$ 中是开集，因此由推论~\ref{cor:open-polish} 是 Polish 空间；若 $A$ 为闭集，则 $f^{-1}(A)$ 在度量空间 $P$ 中闭，而闭集总是 $G_\delta$ 集，因此由定理~\ref{thm:gdelta-polish} 仍是 Polish 空间；若 $A$ 本身就是 $G_\delta$ 集，则 $f^{-1}(A)$ 同样是 $G_\delta$ 集，仍由定理~\ref{thm:gdelta-polish} 得到 Polish 性。将 $f$ 限制到 $f^{-1}(A)$，得到连续双射
\[
f|_{f^{-1}(A)}\colon f^{-1}(A)\to A.
\]
故 $A$ 为 Lusin 空间。
\end{proof}

\paragraph{关于 Suslin 定理与 Lusin 分离定理}

\begin{remark}
完整的描述集合论会进一步研究解析集、共解析集以及 Borel 结构的判别，例如著名的 Suslin 定理和 Lusin 分离定理。对本书而言，真正需要的是前述层级关系与保持性：它们已经足以支撑第 2 章的拓扑测度论、第 6 章以后的可测对偶构造，以及第 11 章中的策略空间建模。
\end{remark}

于是，第 \ref{chap:01} 章到这里也就闭环了。第 \ref{sec:1-1} 节告诉我们极限不能只靠序列，第 \ref{sec:1-2} 节告诉我们存在性需要完备性与 Baire 方法，第 \ref{sec:1-3} 节给出 Polish 这一默认舞台，而本节则说明为什么还要允许 Polish 的连续像。下一章进入测度论与随机空间时，这些工具就不再像抽象预备，而会直接变成优化对象的合法宿主。

\begin{example}[最小抽象例子：连续像出现的状态空间]\label{ex:continuous-image-space}
设 $P$ 为 Polish 空间，$f\colon P\to X$ 为连续满射。即便 $X$ 不能自然嵌入某个线性空间，只要它仍是这样一个连续像，我们就可以把可测结构和逼近论证先放在 $P$ 上做，再通过 $f$ 推到 $X$ 上。把这个例子放在这里，是为了说明为什么定义~\ref{def:suslin-space} 对随机控制和可测选择如此关键。
\end{example}

\begin{example}[有限维坍缩：紧致度量空间与欧氏紧集]\label{ex:compact-metric-lusin}
每个紧致度量空间都是 Polish 空间，因此由命题~\ref{prop:polish-lusin-suslin} 自动也是 Lusin 与 Suslin 空间。特别地，$\mathbb R^n$ 的任意紧子集都处在这一层级中。把这个例子放在这里，是为了说明有限维优化里出现的紧可行域，完全不会触发新的描述集合论困难。
\end{example}

\begin{example}[应用触点：策略与观测的可数编码]\label{ex:policy-parameterization}
在离散时间随机控制中，策略往往由一串历史观测映射决定。若把所有可行观测轨道组织成 Polish 空间，再把策略参数写成标准 Baire 空间上的编码，就可借助命题~\ref{prop:baire-space-surjection} 和命题~\ref{prop:suslin-lusin-stability}，把策略集视为 Suslin 或 Lusin 空间。把这个例子放在这里，是为了给后文第 11 章中的可测参数化问题准备舞台。
\end{example}

\subsection*{有限维对照}

在 Boyd 的有限维语境中，状态空间通常就是 $\mathbb R^n$、其开子集或闭子集，因此完全没有必要再区分定义~\ref{def:polish-space}、定义~\ref{def:lusin-space} 和定义~\ref{def:suslin-space}。但一旦模型涉及连续像、商空间、策略集或路径空间，这三层结构就会开始分家，而命题~\ref{prop:polish-lusin-suslin} 提供了后文回链时最常用的导航。

\subsection*{习题（待补）}

\begin{exercise}[标准 Baire 空间占位]\label{exr:baire-space-placeholder}
补一题：要求将某个具体 Polish 空间写成标准 Baire 空间的连续像，并说明这一表示在可测参数化中的意义。
\end{exercise}

\begin{exercise}[Lusin 子空间占位]\label{exr:lusin-placeholder}
补一题：设 $X$ 为 Lusin 空间，证明其某个给定的 $G_\delta$ 子集仍为 Lusin 空间，并明确指出定理~\ref{thm:gdelta-polish} 在证明中的使用位置。
\end{exercise}

